% docs/manual/developer/chapters/09-gui-extension.tex
% 第 9 章：GUI 扩展

\chapter{GUI 扩展}

\modname{gui} 子包基于 \modname{PySide6}（Qt 6 Python 绑定），实现 14 步向导式配置生成器 + 运行监控 + 远程提交。本章给开发者：架构总览、关键文件、添加新 wizard page 的 walkthrough。

\warninline{GUI 在无图形系统的 HPC 节点上无法启动。开发 GUI 时需要本地 macOS / Linux Desktop / Windows 环境。}

\section{文件结构}

\file{src/openbench/gui/}：

\begin{itemize}
  \item \modname{app}：应用 bootstrap，QApplication 启动
  \item \modname{main\_window}：主窗口（侧栏 + 中心 wizard 区）
  \item \modname{controller}：wizard 控制器（页面切换、状态聚合）
  \item \modname{config\_manager}：schema dataclass ↔ GUI 状态双向同步
  \item \modname{validation} / \modname{data\_validator}：实时校验
  \item \modname{runner}：本地评估运行器（包装 \modname{runner.local}）
  \item \modname{remote\_runner}：SSH 远程运行
  \item \modname{progress\_parser}：把 \cli{run} 输出解析为 GUI 进度事件
  \item \modname{path\_utils}：路径工具
  \item \modname{wizard\_config}：wizard 配置
  \item \modname{pages/}：14 个 wizard page + base\_page
  \item \modname{widgets/}：复用部件（数据源编辑器、模型定义编辑器、远程配置等）
  \item \modname{dialogs/}：弹窗（scan\_confirm 等）
  \item \modname{styles/}：QSS 样式表
\end{itemize}

\section{Qt 架构基础}

\openbench{} GUI 用经典 Qt 模式：

\begin{itemize}
  \item \textbf{Signal/Slot}：组件通信。每个 page 暴露 \texttt{configChanged} 信号；controller 连接到 \texttt{config\_manager} 的 update slot
  \item \textbf{Model/View}（部分 page 用）：QAbstractItemModel + QTableView
  \item \textbf{Async via QThread}：长任务（evaluation run）在子线程，主线程发 UI 事件
\end{itemize}

\subsection{开发用 PySide6 而非 PyQt}

PySide6 是 Qt 公司官方 Python 绑定（LGPL）。PyQt 是 Riverbank 商用替代。两者 API 99\% 兼容。\openbench{} 选 PySide6 是因为 LGPL 友好。

\section{14 个 wizard page}

\file{gui/pages/}：

\begin{longtable}{l p{4cm} p{5cm}}
  \toprule
  page 文件 & schema 对应 & 主要 widget \\
  \midrule
  \endhead
  \modname{page\_general} & project 标识 + 时空 & 文本/数字输入 \\
  \modname{page\_variables} & evaluation.variables & checkable tree（按 category） \\
  \modname{page\_ref\_data} & reference.sources & per-variable 下拉 \\
  \modname{page\_sim\_data} & simulation 字典 & 表格 + entry 编辑器（widgets/data\_source\_editor） \\
  \modname{page\_evaluation} & project 评估开关 & checkbox 组 \\
  \modname{page\_metrics} & metrics list & checkable list \\
  \modname{page\_scores} & scores list & checkable list \\
  \modname{page\_comparisons} & comparison.\{enabled,items\} & checkbox + list \\
  \modname{page\_statistics} & statistics.\{enabled,items\} & 同上 \\
  \modname{page\_runtime} & project.\{num\_cores,force,debug\_mode,...\} & slider + checkbox \\
  \modname{page\_registry} & 浏览/编辑 catalog & 表格 + register 表单 \\
  \modname{page\_options} & project 高级 & 杂项输入 \\
  \modname{page\_preview} & 综合 yaml 预览 + check & 文本 + 状态 \\
  \modname{page\_run\_monitor} & 运行进度 + 日志 & 进度条 + 日志窗 \\
  \bottomrule
\end{longtable}

\section{base\_page：所有 page 的基类}

\file{gui/pages/base\_page.py}：

\begin{minted}{python}
class BasePage(QWidget):
    """所有 wizard page 的基类。"""

    configChanged = Signal()  # 子类发出此信号

    def __init__(self, controller, parent=None):
        super().__init__(parent)
        self.controller = controller
        self._setup_ui()

    def _setup_ui(self):
        """子类 override 创建 widget 树。"""
        raise NotImplementedError

    def load_config(self, cfg: OpenBenchConfig):
        """从 schema 加载到 widget。"""
        raise NotImplementedError

    def save_to_config(self, cfg: OpenBenchConfig):
        """从 widget 写回 schema。"""
        raise NotImplementedError

    def validate(self) -> list[str]:
        """返回错误消息列表（空表示通过）。"""
        return []
\end{minted}

\section{config\_manager：双向同步}

\file{gui/config\_manager.py} 是 schema ↔ widget 状态的中介：

\begin{itemize}
  \item 持有当前 \texttt{OpenBenchConfig} 实例
  \item 监听各 page 的 \texttt{configChanged}，调每个 page 的 \texttt{save\_to\_config}
  \item 提供 \texttt{load\_yaml(path)} / \texttt{save\_yaml(path)}
  \item Preview page 通过它拿"含默认值的完整 yaml 字符串"
\end{itemize}

\subsection{何时同步}

\begin{itemize}
  \item 用户切 page（next / back）→ 立刻保存当前 page 状态
  \item 编辑某 widget → 发 \texttt{configChanged}，但不立刻保存（避免每个键击都序列化）
  \item File → Save → 强制全部 page 保存 + 写盘
\end{itemize}

\section{progress\_parser：CLI 输出 → GUI 事件}

\openbench{} GUI 的 Run Monitor 不直接调 Python API（避免阻塞主线程），而是 fork \cli{openbench run} 子进程，通过管道读 stderr。

\modname{progress\_parser} 解析每行：

\begin{itemize}
  \item \texttt{Phase: ...} → 切换阶段标签
  \item \texttt{Queued <var>: sim=X ref=Y} → 增加任务计数
  \item \texttt{Cache hit / Output written} → 更新进度条
  \item \texttt{ERROR / WARN} → 在日志窗口高亮
\end{itemize}

\noteinline{这种基于文本解析的设计有缺点：CLI 输出格式变了 GUI 进度可能错。但解耦换来稳定性——GUI 不依赖 \openbench{} 内部 API。}

\section{Walkthrough：添加新 wizard page}

假设要加一个"Performance Benchmark page"（让用户预估评估耗时）：

\subsection{Step 1：建文件 \file{gui/pages/page\_performance.py}}

\begin{minted}{python}
from PySide6.QtWidgets import QVBoxLayout, QLabel, QPushButton
from .base_page import BasePage


class PerformancePage(BasePage):
    title = "Performance"

    def _setup_ui(self):
        layout = QVBoxLayout(self)
        self.estimate_label = QLabel("Estimated runtime: --")
        self.run_estimate_btn = QPushButton("Estimate")
        self.run_estimate_btn.clicked.connect(self._estimate)
        layout.addWidget(self.estimate_label)
        layout.addWidget(self.run_estimate_btn)

    def _estimate(self):
        cfg = self.controller.current_config
        # 简单估算：variables × sims × refs × time_factor
        n_var = len(cfg.evaluation.variables)
        n_sim = len(cfg.simulation)
        n_yr = cfg.project.years[1] - cfg.project.years[0] + 1
        seconds = n_var * n_sim * n_yr * 8
        self.estimate_label.setText(f"Estimated runtime: ~{seconds // 60} min")

    def load_config(self, cfg):
        pass  # 此 page 无持久状态

    def save_to_config(self, cfg):
        pass

    def validate(self):
        return []
\end{minted}

\subsection{Step 2：注册到 main\_window}

\file{gui/main\_window.py}：

\begin{minted}{python}
from .pages.page_performance import PerformancePage

# 在 _setup_pages() 中：
self.pages.append(PerformancePage(self.controller))
\end{minted}

\subsection{Step 3：调整 page 顺序}

按用户工作流插入到合理位置（例如 runtime 之后、preview 之前）。

\subsection{Step 4：测试}

GUI 测试目前主要靠人工。建议至少：

\begin{itemize}
  \item 启动 GUI 看新 page 出现
  \item 切到该 page 不报错
  \item 点 Estimate 按钮看显示更新
\end{itemize}

如要自动化，参考 \href{https://pytest-qt.readthedocs.io}{pytest-qt}（headless 测试 Qt 应用）。

\subsection{Step 5：文档}

加到用户卷第 8 章 GUI page 列表。

\section{GUI 测试现状}

\openbench{} 的 GUI 部分目前没有自动测试（headless 环境难以渲染 Qt）。开发 GUI 时主要靠：

\begin{itemize}
  \item 手动测试启动 + 切 page + 编辑 + 保存
  \item Qt 内部 layout 用 QSS 调整时可视化检查
  \item 进度条 / 日志解析的 \modname{progress\_parser} 可以离线 unit test
\end{itemize}

未来想加 pytest-qt 自动测试，欢迎贡献。

\section{下一步}

下一章详解测试与 CI。
